\scanpage{5}
\uline{Corollaire 2} Supposons qu'on ait deux
$F_1$ \unsure{resp$^t$} $F_2$ \inserted{q. coh.} sur $X_1$ \inserted{\unsure{resp|resp.} $X_2$}, et qu'on considère les
\deleted{faisceaux}\inserted{Modules} quotients de $F_1$, $F_2$, $F_1\otimes_V F_2$. Alors
la correspondance du Cor. 1 est compatible
avec la formation du \deleted{\unsure{pro} \unsure{f}} produit tensoriel
de quotients des \deleted{\unsure{$F_i$|$F_1$}, \unsure{faisceaux} \unsure{quotients}}\inserted{$F_1$ \deleted{\unsure{et}} et $F_2$} ($F'_1$ et $F'_2$)
identifié à un quotient de $F_1\otimes F_2$ ($F'_1\otimes F'_2\unsure{\neq|{+}}$)

En effet, le produit tensoriel \uline{sur $V$}
de faisceaux qu. coh. $V$-plats est $V$-plat,
et la \deleted{\unsure{pass}} formation du produit tensoriel
commute à la tensorisation par $K$.

\uline{\deleted{\unsure{P}}Corollaire 3} \deleted{Supposons que} Les \unsure{souspréschémas|sous-préschémas}
fermés $Y$ de $X$ plats sur $\unsure{N|V}$ sont en
correspondance biunivoque avec les
sous-préschémas fermés $Y'$ de $X'=X\otimes_V K$,
à $Y$ correspond $Y'=Y\otimes_V K=Y\cap X'$, et
à $Y'$ \unsure{correspond|correspondent} le sous-préschéma
adhérence $Y=\overline{Y'}$. Cette correspondance
est compatible avec les produits
cartésiens sur $V$.
% archive sheet number: 4

